% Generated draft: requires mathematical and source review.
\documentclass[11pt]{article}
\usepackage[margin=25mm]{geometry}
\usepackage{fontspec}
\setmainfont{Latin Modern Roman}
\usepackage{amsmath,amssymb,mathtools,unicode-math}
\setmathfont{Latin Modern Math}
\usepackage{enumitem,needspace,xurl,hyperref}
\setlength{\parindent}{0pt}
\setlength{\parskip}{5pt}
\setlength{\emergencystretch}{3em}
\newcommand{\N}{\mathbb{N}}
\newcommand{\Z}{\mathbb{Z}}
\newcommand{\Q}{\mathbb{Q}}
\newcommand{\R}{\mathbb{R}}
\newcommand{\C}{\mathbb{C}}
\newcommand{\F}{\mathbb{F}}
\newcommand{\A}{\mathbb{A}}
\newcommand{\PP}{\mathbb{P}}
\newcommand{\E}{\mathbb{E}}
\newcommand{\eps}{\varepsilon}
\newcommand{\dd}{\,\mathrm{d}}
\begin{document}
\section*{UnsolvedMath ID: 30004109}
\textbf{Source problem code:} catalogue entry
\subsection{Definitions and mathematical statement}
The supplied record concerns a block idempotent \(b\) of positive defect in the group algebra \(\mathcal{O}G\). Here the intended setting involves a finite group \(G\), a prime \(p\), and coefficient objects \(\mathcal{O}\) and \(K\) for modular character theory; the record does not specify the hypotheses on these coefficient objects or their relationship. A block idempotent means a primitive central idempotent of the relevant group algebra. The convention for its defect is not supplied.

The notation \([\mathcal{P}(G)/G]\) is described as indexing conjugacy classes of chains of nontrivial \(p\)-subgroups of \(G\). The record does not define the permitted chains, whether the empty chain is included, or the convention for their length \(|\sigma|\). It also does not define \(G_\sigma\) or the construction of \(b_\sigma\), beyond referring to corresponding Brauer-correspondent blocks. Consequently these symbols retain their source meanings pending clarification.

For a parameter tuple \((d,r,L,h)\), the notation
\[
\operatorname{Irr}_K(G_\sigma,b_\sigma,d,r,L,h)
\]
is described in the supplied record as the set of irreducible \(K\)-characters in those corresponding blocks with specified defect \(d\), residue \(r\), character field \(L\), and Hasse invariant \(h\). The admissible parameter domains and the definitions and normalizations of these invariants are not supplied. In particular, the meaning of irreducibility over \(K\) and the required splitting-field assumptions are unresolved.

\textbf{Source conjecture, with unresolved conventions.} For every positive-defect block idempotent \(b\) of \(\mathcal{O}G\) and every admissible tuple \((d,r,L,h)\), the proposed identity is
\[
\sum_{\sigma\in[\mathcal{P}(G)/G]}
(-1)^{|\sigma|}
\left|\operatorname{Irr}_K(G_\sigma,b_\sigma,d,r,L,h)\right|
=0.
\]
This records the supplied mathematical target. A fully specified conjecture requires the missing coefficient hypotheses, chain conventions, block constructions, and parameter definitions; no choices for them are asserted here.
\subsection{Short English statement}
The proposed cancellation identity says that, for each positive-defect block and each admissible choice of four character invariants, the alternating count of the specified characters over conjugacy classes of subgroup chains is zero. The supplied record leaves essential conventions undefined, so this is a qualified statement-only draft requiring mathematical review.
\subsection{Sources}
\begin{itemize}
\item Local UnsolvedMath registry, permanent ID \(30004109\), title \emph{Character-Counting Chain Cancellation in Blocks}, source problem code \texttt{catalogue entry}, collection \emph{Oberwolfach Reports Open Problems}. The supplied record lists proposed year \(2019\) and names no proposer.
\item Supplied source citation: \emph{Representations of Finite Groups (2020)}, Oberwolfach Reports, pp.~841--895, DOI \texttt{10.4171/owr/2019/14}. Supplied source URL and literature-source link, labeled ``Oberwolfach representation theory report'': \url{https://doi.org/10.4171/owr/2019/14}. The background describes these reports as documenting mathematical workshops and their open-problem discussions. No specific contribution, page, or numbered conjecture within the report is identified.
\item Literature assessment dated \texttt{2026-08-22}: the supplied curation labels the problem partially solved and describes it as an extracted Turull-special-case conjecture, with proved cases and structural reductions reported in primary literature and the general conjecture unresolved. This is supplied dated curation, not an independently verified literature assessment. The background expressly requests independent verification of time-sensitive status and cited claims.
\item The supplied registry labels the statement \texttt{corrected\_verified} and reports manual semantic, typographic, and LaTeX review. These are source metadata, not independent confirmation here that the formulation is complete or mathematically correct.
\end{itemize}
\end{document}
